% Input fragment; uses the nonzero-anchor lemma from the 854 proof.
\subsection{Combining the anchor with the Pell equation: 853 variables}
\label{subsec:explicit-853}

One more coordinate can be removed from the explicit construction. The
two algebraic choices to be combined are the private quadratic anchor
and the monic Pell equation for the middle proxy. First, however, one
must remove denominator guards which depend on that proxy. Their removal
is justified over $\mathbb Q$ by the equations themselves.

\begin{lemma}\label{lem:pole-denominator-forced}
Let $b,t,e,p,q,s\in\mathbb Q$, put $B=17b$ and $\delta=1+bt^2$, and set
\[
 K=B\delta^2-b^2(1-bt^2)^2.
\]
Suppose
\begin{align*}
 p^2-B(1-e)(1+3e)&=0,\\
 Bq^2\delta^4-(s^2\delta^2-BeK)(s^2\delta^2+3BeK)&=0.
\end{align*}
Then $\delta\ne0$.
\end{lemma}
\begin{proof}
If $\delta=0$, then $t\ne0$ and $b=-1/t^2$. Substituting this in
the second equation gives
\[
 \frac{48\cdot17^2e^2}{t^{12}}=0,
\]
so $e=0$. The first equation, multiplied by $t^2$, now becomes
$(pt)^2+17=0$, which has no rational solution.
\end{proof}

For $\delta\ne0$, set
\[
 U_{\rm pole}=\frac{b(1-bt^2)}\delta,
 \qquad h=e(B-U_{\rm pole}^2)=\frac{eK}{\delta^2}.
\]
The second equation of the lemma is exactly $\delta^4$ times
\[
 Bq^2-(s^2-Bh)(s^2+3Bh)=0.
\]
Thus the two cleared conics already imply the excluded denominator is
nonzero. This applies to all $42$ pole blocks. If $b$ itself is a rational
expression, its denominator and all other denominator factors are kept
nonzero; multiplication by those nonzero factors does not affect the
argument. In the actual circuit, each second conic has its own omitted
denominator numerator to the fourth power, and no other omitted factor.
The other $255$ residual denominators use only retained factors.

After deleting these $42$ guards, there are $1639$ guard numerators,
with $1560$ distinct nonconstant factors. Let $G$ be the product of
those factors. Exact inspection of the circuit proves that $G$ depends
only on the input, the $633$ point coordinates, the coordinate
denominator, and the two contact coordinates. In particular it is
independent of all proxy and anchor coordinates. The guard deletion is
an equivalence of rational solution sets. It is not a claim about
geometric fibers over an algebraic closure, where the last step of
Lemma~\ref{lem:pole-denominator-forced} would fail.

\begin{lemma}\label{lem:two-quadratics-one-coordinate}
Let $k$ be a field of characteristic different from two, and let
$N,\Delta\in k$. Define
\[
 h(W)=W^4+2(N-\Delta)W^2+(N+\Delta)^2.
\]
For every nonzero $W\in k$, the equation $h(W)=0$ is equivalent to
the existence of $T,V\in k$ such that
\[
 T^2=-N,\qquad V^2=\Delta,\qquad W=T+V.
\]
The corresponding coordinates are uniquely given by
\[
 T=\frac{W^2-N-\Delta}{2W},\qquad
 V=\frac{W^2+N+\Delta}{2W}.
\]
\end{lemma}
\begin{proof}
If $W=T+V$, then $(V-T)W=V^2-T^2=N+\Delta$, which gives the
displayed reconstruction formulas. Conversely those formulas give
$T+V=W$. Direct expansion gives both identities
\[
 4W^2(T^2+N)=h(W),\qquad
 4W^2(V^2-\Delta)=h(W).
\]
Since $2W\ne0$, these establish the equivalence.
\end{proof}

\paragraph{Application to the anchor and proxy.}
As in the preceding construction, the cleared anchor is $Bu^2+N=0$
with $B=S^2$. Its three witnesses are private, and
Lemma~\ref{lem:nonzero-anchor-witness} allows $u\ne0$. The retained
guards imply $S\ne0$. Put $T=Su$, so $T^2=-N$ and $T\ne0$.

Write the middle proxy as $\tau$. Its equation is
\[
 \tau^2-a\tau+b_0=0,
 \qquad a=s(s-2)z,
 \qquad b_0=(s-2)^2z^2-1.
\]
Set $\Delta=a^2-4b_0$ and $V=2\tau-a$. Then the proxy equation is
equivalent to $V^2=\Delta$. The sign of $T$ is free: changing it
changes only the sign of the private anchor witness. Because $T\ne0$,
at least one of $T+V$ and $-T+V$ is nonzero. Thus one may use
$W=T+V\ne0$ after choosing that sign. No additional finite choice
of a primitive-element coefficient is required.

Put $W=1/(GU)$ and define
\begin{align*}
 \Psi&=1+2(N-\Delta)G^2U^2+(N+\Delta)^2G^4U^4,\\
 P_\tau&=1+2aGU+(N+\Delta)G^2U^2,\\
 Q_\tau&=4GU.
\end{align*}
The coefficients $N,\Delta,G,a$ are independent of the new witness
$U$. Hence $\Psi=1$ whenever $GU=0$. For $GU\ne0$,
\[
 \Psi=(GU)^4 h\bigl(1/(GU)\bigr),\qquad
 \tau=\frac{P_\tau}{Q_\tau}.
\]
The latter formula follows from $\tau=(a+V)/2$ and
Lemma~\ref{lem:two-quadratics-one-coordinate}. Explicitly,
\[
 T=\frac{1-(N+\Delta)G^2U^2}{2GU},\qquad
 V=\frac{1+(N+\Delta)G^2U^2}{2GU}.
\]
The recovered anchor witness is $u=T/S$. It may be zero in this
reverse construction, which is allowed by the original equations.

\paragraph{Substitution without additional variables.}
Remove the old anchor and proxy equations, leaving $295$ cleared
residual numerators $F_i$. For each choose a syntactic degree bound
$d_i$ for its degree in $\tau$, and define
\[
 \widetilde F_i
  =Q_\tau^{d_i}
      F_i\bigl(\tau=P_\tau/Q_\tau\bigr).
\]
This is a polynomial. It is constructed by an arithmetic circuit,
with no auxiliary quantified variables. At an addition gate, pad
the lower denominator power to the maximum of the children's powers;
at a multiplication gate, add the powers; at a power gate, multiply
the power. This proves the displayed substitution identity
inductively, even when the chosen syntactic degree is larger than
the actual degree after cancellation.

For each fixed branch, form
\[
 Q=\sum_{i=1}^{295}\widetilde F_i^2+\Psi^2.
\]
Every rational solution of the original system first admits the
nonzero-anchor normalization, then a choice of sign with $W\ne0$,
and finally $U=1/(GW)$. The reconstructed proxy is the original
$\tau$, so every substituted residual vanishes.

Conversely, $Q=0$ forces $\Psi=0$ and all
$\widetilde F_i=0$. Since $\Psi=1$ at $GU=0$, both $G$ and
$Q_\tau$ are nonzero. The reconstruction identities recover the
anchor and proxy equations. The substitution identities then give
$F_i=0$ for every old residual. The retained guards hold, and
Lemma~\ref{lem:pole-denominator-forced} recovers the $42$ omitted
guards. Consequently all old rational equations are valid and the
old independent inverse witness can also be restored. The input
projection is therefore unchanged.

Take the product of $Q$ over the same $212\cdot256=54272$ finite
chart and height choices as before. This is one integer polynomial
in $853$ variables, including the input, with $852$ existential
witnesses. The old inverse coordinate and middle proxy coordinate
have been removed, and the old first anchor coordinate has been
reused as $U$. The proof permits changes in all three private
anchor witnesses and does not assert equality of every fiber
obtained by fixing those witnesses.

\paragraph{Certificate and quantitative bounds.}
The exact specification is
\texttt{polynomial-853/polynomial-circuit.json}. It has $95084$
stored nodes, of which $94946$ are reached from the output, and
all $853$ declared coordinates occur in that output directed
acyclic graph. The greatest homogenizing exponent is $1280$.
The independent reader \texttt{verify\_compact\_853.py} imports
neither producer. It reruns the separate pole-guard verifier,
checks the complete anchor and Pell identities, checks both
quadratic reconstruction identities, and certifies all $295$
substitutions by exact expression-tree induction. It also checks
the coordinate map, finite parameter lists, sum of $296$ squares,
and degree and coefficient-norm bounds. A separate review of the
producer confirms the signs, substitution invariant, and variable
renumbering. These transformation checks use exact arithmetic.

The resulting rigorous bounds are
\[
 \deg P\le894090701148160,
 \qquad H(P)\le2^{29154205340326912}.
\]
They are upper bounds, not exact expansion statistics. The smaller
variable count comes at the cost of a worse degree bound than the
$854$-variable construction. Both specifications are retained.
The arithmetic identification of the input projection with
$\mathbb Z$ still inherits the underlying source theorem; the
transformation certificates do not independently verify every
arithmetic assertion of that theorem.
